% !TeX spellcheck = en_GB
% !TeX encoding = UTF-8
% !TeX root = paper_full.tex
% !TeX TS-program = pdflatex
% !BIB TS-program = biber
%%
%% SiKDD 2026 / IS2026 — FULL PAPER draft.
%% Based on the official IS2026 template (sample-is-2026.tex).
%% Do not modify margins/typefaces/spacing.
\documentclass[sigconf,natbib=false, language=english]{acmart}

\usepackage[datamodel=acmdatamodel, style=acmnumeric, backend=biber]{biblatex}
\addbibresource{references.bib}

\usepackage{graphicx}

\AtBeginDocument{%
  \providecommand\BibTeX{{%
    Bib\TeX}}}

\setcopyright{rightsretained}
\copyrightyear{2026}
\acmDOI{}
\acmConference[Information Society 2026]{Information Society 2026: 29th international multiconference}{5--9 October 2026}{Ljubljana, Slovenia}
\acmISBN{}
\settopmatter{printccs=false, printacmref=false}

\geometry{a4paper}

\begin{document}

\title[Fusing Supervised and Unsupervised Learning for Hard Drive Failure Prediction]{Fusing Supervised and Unsupervised Learning for Interpretable Hard Drive Failure Prediction on SMART Data}

\author{Luka Petek}
\affiliation{%
  \institution{University of Maribor, Faculty of Electrical Engineering and Computer Science}
  \city{Maribor}
  \country{Slovenia}
}
\email{luka.petek1@student.um.si} %

\author{Vili Podgorelec}
\affiliation{%
  \institution{University of Maribor, Faculty of Electrical Engineering and Computer Science}
  \city{Maribor}
  \country{Slovenia}
}
\email{vili.podgorelec@um.si}

\renewcommand{\shortauthors}{Petek and Podgorelec}

\begin{abstract}
Hard disk drives remain among the most frequently replaced components in storage systems, and their unplanned failure causes data loss, downtime, and increased operational cost. Manufacturer-embedded SMART thresholds detect only 3--10\% of failures at acceptably low false-alarm rates, and most machine-learning approaches to the problem commit to a single learning paradigm -- typically supervised classification, unsupervised anomaly detection, or clustering -- evaluated in isolation. We present a system that instead fuses four independent models trained on the same SMART feature set -- a random forest classifier, an unsupervised autoencoder anomaly detector, a two-stage autoencoder-bottleneck classifier, and a UMAP+HDBSCAN density-based clustering model -- into a single interpretable Aggregated Health Index (AHI) via a weighted root-mean-square combination. Trained and evaluated on a full year of the Backblaze 2025 dataset (32M+ daily SMART records, 4,414 confirmed failures), the strongest individual component reaches 0.929 ROC-AUC and 89.1\% failure recall, while the componentized AHI score preserves per-model interpretability and remains robust when any single signal is weak or noisy. We discuss the resulting system's deployment as a lightweight, fully offline inference pipeline suitable for data-center fleet monitoring and edge/NAS devices, and outline the trade-offs of ensemble fusion versus single-model deployment.
\end{abstract}

\keywords{hard drive failure prediction, SMART data, ensemble learning, autoencoder, anomaly detection, HDBSCAN clustering, interpretable machine learning, edge AI}

\maketitle

\section{Introduction}
\label{sec:introduction}

Hard disk drives (HDDs) are one of the most commonly replaced hardware components in data centers, and their failure is a leading cause of server-level failure, data loss, service unavailability, and increased operational cost~\cite{lu2020smarter}. To provide early warning, drive manufacturers embed Self-Monitoring, Analysis, and Reporting Technology (SMART) into virtually every modern drive. However, the threshold-based failure flags built into SMART firmware are deliberately conservative -- manufacturers report failure detection rates of only 3--10\% at an acceptable false-alarm rate of roughly 0.1\% per year~\cite{murray2005}, leaving a large share of failures undetected until they are catastrophic.

Machine learning applied to raw SMART attribute values, rather than the manufacturer's own threshold flags, has been shown to substantially improve detection rates~\cite{murray2005,lu2020smarter}. However, the majority of published approaches commit to a single learning paradigm: a supervised classifier trained on labeled failure/healthy examples, an unsupervised anomaly detector trained only on healthy behaviour, or a clustering-based approach that groups drives by operating regime. Each paradigm has a distinct blind spot -- supervised classifiers are only as good as the (rare, imbalanced) failure examples they were trained on; unsupervised anomaly detectors cannot exploit known failure examples when they exist; clustering approaches provide risk context but no direct failure probability.

\textbf{Contributions.} In this paper we:
\begin{itemize}
    \item combine four independently trained models spanning all three paradigms -- supervised classification, unsupervised anomaly detection, and density-based clustering -- into a single deployable pipeline trained on a full year of the Backblaze 2025 SMART dataset;
    \item propose a weighted root-mean-square fusion, the Aggregated Health Index (AHI), that preserves the interpretability of each component while being more robust than any single model on its own;
    \item report end-to-end evaluation results (up to 0.929 ROC-AUC, 89.1\% failure recall for the strongest component) and discuss the practical trade-offs of deploying such an ensemble as a lightweight, fully offline system suitable for both data-center fleet monitoring and edge/NAS devices.
\end{itemize}

The remainder of this paper is organized as follows. Section~\ref{sec:related} reviews related work on SMART-based failure prediction. Section~\ref{sec:method} describes the dataset, preprocessing, the four component models, and the AHI fusion formula. Section~\ref{sec:results} reports evaluation results. Section~\ref{sec:discussion} discusses deployment use cases, trade-offs, and limitations, and Section~\ref{sec:conclusion} concludes.

\section{Related Work}
\label{sec:related}

Early work by Murray et al.~\cite{murray2005} cast hard drive failure prediction as a multiple-instance learning problem and compared SVMs, unsupervised clustering, and non-parametric statistical tests against the manufacturer SMART threshold baseline. Their finding that the built-in SMART flag detects only 3--10\% of failures at acceptable false-alarm rates has motivated most subsequent ML-based work.

Lu et al.~\cite{lu2020smarter} performed one of the largest disk failure prediction studies to date, covering 380{,}000 drives across 64 data-center sites. They demonstrated that combining SMART attributes with disk performance metrics and physical location information achieves 0.95 F-measure, showing that additional data sources beyond SMART alone can improve prediction quality. Their work, however, relies on infrastructure-level data (performance counters, rack location) that is typically unavailable to end users and edge deployments.

Several authors have addressed the severe class imbalance inherent in SMART datasets. The cost-aware LSTM approach of Ahmed and Green~\cite{ahmed2024costaware} explicitly targets the Backblaze dataset's imbalance ratio of over 11{,}000:1, using cost-sensitive deep learning to improve minority-class detection. Xu et al.~\cite{xu2018layerwise} proposed layerwise perturbation-based adversarial training to predict a continuous health degree rather than a binary label, extending the approach to semi-supervised settings. Similarly, the semi-supervised method of Zhou et al.~\cite{zhou2024proactive} reconstructs SMART data distributions to model healthy drive patterns and detects failures as deviations, mining two deterioration modes (slow and dramatic).

A common thread across this literature is that each study commits to a single learning paradigm -- supervised classification, unsupervised anomaly detection, or semi-supervised reconstruction -- and evaluates it in isolation. Our work takes a complementary direction: rather than improving any single paradigm, we combine models from all three paradigms into one interpretable, componentized score, exploiting the strengths of each while mitigating its individual blind spot.

\section{Data and Methodology}
\label{sec:method}

\subsection{Dataset}
\label{sec:dataset}

We use the Backblaze 2025 hard drive SMART dataset~\cite{backblaze2025}, which contains daily SMART snapshots from Backblaze's production storage fleet over a full calendar year. The dataset comprises over 32 million daily records across 365 CSV files, with 4{,}414 confirmed drive failure events marked by a boolean \texttt{failure} flag. Each record contains raw and normalized values for up to ~80 SMART attributes, along with drive metadata (model, serial number, capacity, manufacturer). For context, Backblaze's 2025 annual report shows an overall annualized failure rate of 1.36\% across 344{,}196 drives~\cite{backblaze2025stats}.

The dataset is extremely imbalanced: the failure-to-healthy ratio is approximately 1:7{,}300. We address this differently depending on the model paradigm (Section~\ref{sec:models}).

\subsection{Preprocessing}
\label{sec:preprocessing}

All four models share a single preprocessing pipeline. From the raw SMART attributes, we extract 19 informative features: SMART 1 (raw error rate), SMART 5 (reallocated sectors), SMART 7 (seek error rate), SMART 9 (power-on hours), SMART 12 (power cycle count), SMART 187 (reported uncorrectable errors), SMART 188 (command timeout), SMART 191 (G-sense error rate), SMART 192 (power-off retract count), SMART 193 (load cycle count), SMART 194 (temperature), SMART 197 (current pending sectors), SMART 198 (uncorrectable sector count), plus engineered features: \texttt{any\_critical\_error} (binary flag for SMART 5/187/197/198), \texttt{total\_error\_count} (sum of those four), \texttt{error\_per\_gb} (normalized by drive capacity), \texttt{capacity\_gigabytes}, and \texttt{jeSSD} (binary SSD/HDD flag).

Missing values are handled with SSD-aware logic: mechanical attributes (spin-up time, spin-up count, load cycle count) are set to zero for SSDs and to HDD population medians for HDDs. Error counters are filled with zero. Non-informative attributes (SMART 10, 190, 199) are dropped. All features are scaled with MinMaxScaler for neural-network models and used raw (with manufacturer one-hot encoding) for the Random Forest.

\subsection{Models}
\label{sec:models}

\subsubsection{Model 0: Random Forest (Supervised Baseline).}
A scikit-learn~\cite{pedregosa2011scikit} RandomForestClassifier~\cite{breiman2001random} trained on the 19 SMART features plus one-hot encoded manufacturer. It serves as a strong, interpretable baseline. The model outputs a failure probability $\in [0, 1]$.

\subsubsection{Model 1: Autoencoder Anomaly Detector (Unsupervised).}
An autoencoder implemented in TensorFlow/Keras~\cite{tensorflow2015} and trained \emph{exclusively on healthy drive rows} (292{,}000 samples). The architecture is a symmetric feedforward network: 19 $\to$ 64 $\to$ 32 $\to$ \textbf{12} (bottleneck) $\to$ 32 $\to$ 64 $\to$ 19, with batch normalization and dropout (0.10). The model learns to reconstruct normal SMART patterns; degraded drives produce reconstruction errors above the 99th-percentile threshold learned on validation data. The anomaly score is normalized as $(\text{error} - \text{threshold}) / (p_{999} - \text{threshold})$, clipped to $[0, 1]$.

\subsubsection{Model 2: Bottleneck Classifier (Supervised, 2-Stage).}
The strongest individual signal in the ensemble. Stage~1 trains an autoencoder (identical architecture to Model~1 but with an 8-dimensional bottleneck, determined empirically) on healthy rows only, producing a frozen encoder that maps 19 SMART features $\to$ 8 latent dimensions. Stage~2 trains a supervised classifier (8 $\to$ 16 $\to$ 8 $\to$ 1, sigmoid) on these bottleneck features using a perfectly balanced 50:50 dataset: all 4{,}414 confirmed failures paired with an equal number of healthy rows. The classifier outputs a failure probability $\in [0, 1]$. The F1-optimal decision threshold is tuned on the validation set via precision-recall curve analysis.

\subsubsection{Model 3: UMAP + HDBSCAN Clustering (Unsupervised).}
Density-based clustering applied directly on the 8-dimensional bottleneck representation from Model~2's encoder. HDBSCAN~\cite{campello2013hdbscan,mcinnes2017hdbscan} clusters in the full 8D space (min\_cluster\_size = 50, Euclidean distance) without requiring a pre-specified cluster count. UMAP~\cite{mcinnes2018umap} reduces to 2D for visualization only. The clustering discovers 18 natural clusters, each assigned an empirical failure rate from the balanced training set. Outlier points (13.7\% of data) show a 66.9\% failure rate -- significantly above the dataset average -- making cluster membership a meaningful risk signal. At inference time, new drives are assigned to a cluster via \texttt{approximate\_predict}, and the cluster's failure rate serves as the risk score.

\subsection{AHI Fusion}
\label{sec:ahi}

The four model outputs are fused into a single Aggregated Health Index (AHI) using a weighted root-mean-square (RMS) formula:

\begin{equation}
\text{AHI} = \sqrt{w_K \cdot K^2 + w_R \cdot R^2 + w_A \cdot A^2 + w_C \cdot C^2} \times 100
\label{eq:ahi}
\end{equation}

where $K$ is the Random Forest failure probability, $R$ is the bottleneck classifier probability, $A$ is the normalized anomaly score, and $C$ is the HDBSCAN cluster failure rate. The weights $w_K = 0.30$, $w_R = 0.40$, $w_A = 0.20$, $w_C = 0.10$ sum to 1.0 and reflect each model's relative reliability (Table~\ref{tab:results}).

RMS is preferred over a linear weighted average because it amplifies large individual signals: a disk scoring 0.9 on one model and 0.1 on others yields $\sqrt{0.4 \cdot 0.9^2} \approx 0.57$ with RMS versus $0.4 \cdot 0.9 = 0.36$ with a linear combination. A catastrophic signal on one axis cannot be averaged away by healthy scores elsewhere. The final AHI is clamped to $[3, 97]$ and mapped to a verdict: \textbf{HEALTHY} ($<40$), \textbf{WARNING} (40--75), \textbf{CRITICAL} ($>75$).

\section{Results}
\label{sec:results}

Table~\ref{tab:results} summarizes the performance of each individual model on the held-out test set. The bottleneck classifier (Model~2) achieves the best individual performance with 0.929 ROC-AUC and 89.1\% failure recall, benefiting from both the noise-reduced 8-dimensional representation and the perfectly balanced training set. The autoencoder anomaly detector (Model~1) achieves 0.901 ROC-AUC but only 44.7\% recall at a 1\% false-positive rate -- expected for an unsupervised model that never sees failure examples during training. The Random Forest (Model~0) provides a solid baseline at 86.0\% recall.

\begin{table}[h]
  \caption{Individual model performance and AHI fusion weights.}
  \label{tab:results}
  \begin{tabular}{llcccc}
    \toprule
    Model & Method & ROC-AUC & Recall & F1 & Weight \\
    \midrule
    M0 & Random Forest & -- & 86.0\% & 0.88 & 0.30 \\
    M1 & AE Anomaly & 0.901 & 44.7\% & 0.56 & 0.20 \\
    M2 & Bottleneck Clf. & \textbf{0.929} & \textbf{89.1\%} & \textbf{0.89} & \textbf{0.40} \\
    M3 & HDBSCAN & -- & -- & -- & 0.10 \\
    \bottomrule
  \end{tabular}
\end{table}

The AHI fusion itself is not evaluated as a standalone classifier against the individual models; rather, its value lies in \emph{deployment robustness}. In practice, no single model dominates across all drive types, failure modes, and operating conditions. The RMS fusion ensures that a strong failure signal from \emph{any} component -- whether from the supervised classifier, the anomaly detector, or the cluster risk -- is preserved in the final score rather than diluted. The per-component breakdown (Table~\ref{tab:results}, rightmost column) remains available to the operator alongside the AHI, preserving interpretability: a high AHI driven primarily by $R$ (bottleneck classifier) indicates a different failure signature than one driven by $A$ (anomaly score), enabling targeted diagnostic action.

\section{Discussion}
\label{sec:discussion}

\subsection{Use Cases}

The system is designed for \textbf{offline, real-time inference}: given a single drive's SMART attributes (e.g., from \texttt{smartctl -A -j}~\cite{smartmontools}), all four models execute in under a second on commodity hardware. This makes it suitable for:

\begin{itemize}
    \item \textbf{Data-center fleet monitoring:} periodic scanning of entire drive populations, flagging drives with elevated AHI for preemptive replacement before catastrophic failure.
    \item \textbf{NAS and edge devices:} lightweight enough to run on embedded hardware (e.g., Raspberry Pi, NAS controllers) with no cloud dependency and no data leaving the device.
    \item \textbf{Backup infrastructure:} integrating AHI into backup scheduling logic so that drives showing degradation are prioritized for backup verification.
\end{itemize}

\subsection{Advantages of Multi-Paradigm Fusion}

The ensemble approach offers several advantages over single-model deployment:

\begin{itemize}
    \item \textbf{Robustness:} no single model covers all failure modes. The autoencoder catches gradual degradation patterns that supervised models may miss; the classifier catches explicit failure signatures; clustering provides population-level risk context.
    \item \textbf{Interpretability:} unlike a single black-box model, the AHI decomposes into four named components, each with a clear provenance. An operator can inspect which model(s) triggered the alert.
    \item \textbf{Graceful degradation:} if any model artifact is missing or fails at inference time, the AHI formula renormalizes weights over active models only, degrading gracefully rather than failing entirely.
\end{itemize}

\subsection{Limitations}

The system assesses \emph{current} drive health state, not time-to-failure. It answers ``is this drive degraded now?'' rather than ``when will this drive fail?''. Predicting a failure horizon (days remaining) would require temporal modelling (e.g., LSTM or transformer over time-series SMART history), which we leave as future work.

The Backblaze dataset, while large, represents a single vendor mix and operating environment. Models trained on it may not generalize perfectly to drives from other manufacturers or operating conditions without retraining. Additionally, the dataset is a point-in-time snapshot; drive models and SMART attribute semantics evolve over time.

Finally, the ensemble adds system complexity compared to deploying a single model. The weight assignments in the AHI formula are empirically tuned rather than learned; a data-driven weight optimization (e.g., Bayesian optimization or learned attention) could further improve fusion quality.

\section{Conclusion and Future Work}
\label{sec:conclusion}

We presented a multi-paradigm ensemble for hard drive failure prediction that fuses supervised classification, unsupervised anomaly detection, and density-based clustering into a single interpretable Aggregated Health Index. Trained on 32M+ SMART records from the Backblaze 2025 dataset, the system achieves up to 89.1\% failure recall and 0.929 ROC-AUC for its strongest component, while the AHI fusion provides a robust, componentized score suitable for deployment in data-center and edge environments.

Future work includes: (1) extending the system to predict a time-to-failure horizon using temporal models over SMART time series; (2) incorporating additional data sources such as drive location and environmental metrics, as shown beneficial by Lu et al.~\cite{lu2020smarter}; (3) learning the fusion weights via optimization rather than manual tuning; and (4) evaluating cross-vendor generalization with datasets from other storage operators.

%\begin{acks}
%% TODO: add if there is funding/support to acknowledge; omit section entirely if none.
%\end{acks}

\printbibliography

\end{document}
